\documentclass[10pt,a4paper]{amsart}
\usepackage[margin=19mm]{geometry}
\usepackage{amsmath,amssymb,mathtools}
\usepackage[scr=rsfs,cal=euler]{mathalfa}
\usepackage{xfrac}
\usepackage[unicode,hidelinks,bookmarks=false]{hyperref}
\newcommand{\ie}{i.e.\ }
\newcommand{\cf}{cf.\ }
\newcommand{\resp}{respectively\ }
\newcommand{\Subsection}{\subsection}
\providecommand{\nobreakdash}{\nobreak\hbox{-}}
\setlength{\emergencystretch}{2em}
\multlinegap0pt
\usepackage{bm}
\usepackage{bbm}
\usepackage{dsfont}
\usepackage{xspace}
\usepackage{cancel}
\usepackage{mathtools}
\usepackage{amsthm}
\makeatletter
\@ifundefined{subsection}{\newtheorem{subsection}{Subsection}}{}
\@ifundefined{proposition}{\newtheorem{proposition}[subsection]{Proposition}}{}
\makeatother
\newcommand{\Aut}{\operatorname{Aut}}
\newcommand{\Q}{{\mathbb Q}}
\newcommand{\V}[2]{V_{(#1,#2)}}
\newcommand{\eig}{{\mathrm{eig}}}
\newcommand{\kg}{{\mathfrak g}}
\title[Non-existence of Enriques manifolds from OG10 type manifolds]{Non-existence of Enriques manifolds from OG10 type manifolds}
\author{Simone Billi, Franco Giovenzana, Luca Giovenzana, Annalisa Grossi}
\date{Source: arXiv:2501.06893v2 (CC BY 4.0)}
\begin{document}
\maketitle

\noindent\emph{Source and licence.} Non-existence of Enriques manifolds from OG10 type manifolds. Version arXiv:2501.06893v2 (CC BY 4.0), distributed under the Creative Commons Attribution 4.0 International licence, as recorded in the arXiv licence metadata for this exact version and shown on the abstract page https://arxiv.org/abs/2501.06893v2. Full rights notice: licenses/arxiv-2501.06893-CC-BY-4.0.txt.\par\smallskip

\noindent\textit{Editorial scope.} Counted unit: the Proposition whose source label is \texttt{prop:AB-eig} in the article (source0.tex lines 327--332, source label \texttt{prop:AB-eig}), delivered together with the article's own complete original proof of it (source0.tex lines 333--363, one \texttt{proof} environment of 31 lines). No mathematics is written, repaired or re-derived: statement and proof are the article's own text line for line. Required context retained uncounted: decomposition (lines 314--316). Every definition, display and label of those lines is retained unchanged. Omitted from this package and not redistributed: 1--313, 317--326, 364--636. Nothing inside a retained excerpt is omitted or altered: every retained body is delivered exactly as the article's lines are written, so no omission receipt of any kind accompanies this package. Citations: the retained excerpts carry no citation command at all, so no bibliography entry of the article is transcribed and no reference is redistributed. Reference adaptation: none was needed and none was made; every reference inside the delivered unit resolves inside it, so no unresolved reference or citation is produced. Printed slips kept literal: no printed word, symbol, hypothesis or number of the counted statement or of its proof was altered. Layout and class: the article's own class is \texttt{amsart} with packages including anysize, geometry, comment, xcolor, amsmath, mathtools, xy, inputenc, varioref, ulem, amsfonts, amssymb, bbm, esint; the wrapper is the lane's pinned \texttt{amsart} editorial wrapper at 10pt on A4, which restates the theorem-like environments the retained text needs and adds the article's own macro definitions that the retained text needs (\texttt{\textbackslash Aut}, \texttt{\textbackslash Q}, \texttt{\textbackslash V}, \texttt{\textbackslash eig}, \texttt{\textbackslash kg}), with extra wrapper packages (bm, bbm, dsfont, xspace, cancel, mathtools). The delivered numbering is the unit's own. Assets: no retained span contains a figure environment, a graphics-inclusion command, a PDF-page inclusion, a table or a TikZ picture; no image file is copied into this package, the package declares no asset dependency and no image file is redistributed with it. Licence: version arXiv:2501.06893v2 of this article is distributed under the Creative Commons Attribution 4.0 International licence, as recorded in the arXiv licence metadata for this exact version and shown on the abstract page https://arxiv.org/abs/2501.06893v2. A full rights notice is delivered at licenses/arxiv-2501.06893-CC-BY-4.0.txt. Characters: every retained excerpt is pure ASCII, so the delivered engine, XeLaTeX with its default Latin Modern text font, reports no missing character.
        \begin{align}\label{decomposition}
            H^*(X,\mathbb{Q}) = A \oplus B
        \end{align}

    \begin{proposition}\label{prop:AB-eig}
        Let $\varphi\in \Aut (X)$ be of finite order. Then $\varphi^*$ preserves the LLV decomposition \eqref{decomposition}, i.e. $\varphi^*(A) = A$ and $\varphi^*(B) = B$, so that
        \[
        \eig (\varphi^*, H^*(X,\Q),1) = \eig (\varphi^*, A,1) + \eig(\varphi^*, B,1).
        \]
    \end{proposition}

    \begin{proof}
        As $A= \mathfrak g \cdot H^0$ we can write any element $x\in A$ as
        \[
        x = \beta + \sum \alpha_{i_1} \cup \cdots\cup \alpha_{i_r}
        \]
        for some $\beta\in H^0$ and $\alpha_{i_j}\in H^2$, so that we compute 
        \[
        \varphi^*x = \varphi^*\beta + \sum \varphi^*\alpha_{i_1} \cup \cdots\cup \varphi^*\alpha_{i_r}
        \]
        and conclude that $\varphi^* x\in A$, as wanted.

        Now we turn our attention to the $\kg$-submodule $B\simeq\V22$. First, notice that the direct sum
        \[
        H^{10}(X,\Q) = \left( A \cap H^{10}(X,\Q) \right) \oplus \left( B \cap H^{10}(X,\Q)\right)
        \]
        is orthogonal with respect to the intersection form. Indeed, the intersection product of an element $a\in A\cap H^{10}$ with an element $b\in B\cap H^{10}$ can be expressed in terms of the $\kg$-action on $b$, and \(B\cap H^{20}=0\).
        As $\varphi^*$ acts as an isometry for the intersection form and $\varphi^* (A) = A$, for any $x\in B \cap H^{10}$ we have $\varphi^* x\in B \cap H^{10}$. As $B\simeq\V22$ is irreducible, we can write $B = \mathfrak g \cdot x$ for any nonzero element $x\in B \cap H^{10}$, so that any element $y\in B$ can be written as finite sum of elements of the form:
\begin{align*}
    g_1\cdot g_2\cdots g_r \cdot x
\end{align*}
where $g_i=L_{\alpha}$ or $g_i=\Lambda_{\beta}$ for some $\alpha, \beta \in H^2(X,\mathbb Q)$.
        We compute
        \[
        \varphi^* (L_\alpha\cdot x) = \varphi^* (\alpha\cup x) = \varphi^* \alpha\cup \varphi^* x = L_{\varphi^* \alpha}\cdot \varphi^* x
        \] 
        and similarly for the dual Lefschetz operator $\Lambda_\beta := *^{-1}\circ L_\beta \circ *$ for some $\beta\in H^2$, i.e. $\varphi^*(\Lambda_\beta \cdot x) = \Lambda_{\varphi^*\beta} \cdot \varphi^* x$. We conclude that $\varphi^* y\in B$ as it is the finite sum of 
\begin{align*}
    \varphi^* (g_1\cdot g_2\cdots g_r \cdot x ) = g_1'\cdot g_2'\cdots g_r' \cdot \varphi^*x
\end{align*}
where $g_i'\in \kg$ and $\varphi^*x \in B$ as discussed above.
    \end{proof}

\end{document}
